\section{Productive Preservation and the Frontier--Closure Quotient}
\label{sec:or-productive-equivalence}

For a finite library $A$ containing the inactive strategy, write
\[
 F_A(b)=\max_{s\in A}j_s(b),\qquad
 C_A=\cl\!\left(\bigcup_{s\in A}\mods(s)\right),\qquad
 K_A=(F_A,C_A).
\]

The forward preservation result uses only the process factorization and local
update assumptions stated in Section~S1. A converse needs a distinct
detectability assumption: whenever two libraries have equal frontiers but
unequal closures, at least one availability-tagged project cost, duration,
incumbent-reward block, or joint terminal observation differs after the
declared generation, verification, admission, and projection steps.

\begin{theorem}[Frontier--closure characterization]
\label{thm:or-frontier-closure-characterization}
If $F_A=F_L$ and $C_A=C_L$, then $A$ and $L$ induce identical
productive observations and identical finite-horizon values and maximizing
actions. The stationary conclusions of Section~S1 also agree under its
discount clause. If closure is detectable in the preceding sense, equality
of all productive observations together with frontier equality implies
$C_A=C_L$.
\end{theorem}

\begin{proof}
Frontier--closure equality is $K_A=K_L$. The factorization assumption then
identifies project availability, cost, duration, incumbent operating rewards,
generation and verification laws, and the conditional joint completion law.
The normalized admission map and deterministic pushforward from Section~S1
therefore identify the complete productive transition observations. The
finite-horizon and stationary value and action claims follow from
Theorem~\ref{thm:or-raw-compressed-projection}. Conversely, suppose the
frontiers and all named productive observations agree. If the closures
differed, detectability would produce a named observation that differs, a
contradiction. Hence the closures agree. Detectability is used only for this
reverse implication.
\end{proof}

\begin{corollary}[Source-relative safe compression]
\label{cor:or-safe-compression}
For $s_0\in A\subseteq L$, exact equality $F_A=F_L$ and $C_A=C_L$
is sufficient for productive preservation. With exact positive active
weights and $w_{s_0}=0$, the finite safe-feasible family contains $L$
and therefore has at least one minimum-burden member.
\end{corollary}

\begin{proof}
The productive claim is the forward direction of the theorem. The source is
safe relative to itself. A finite source has finitely many sublibraries, so
the nonempty feasible set has an attained minimum under exact weights.
\end{proof}

The converse characterization is deliberately restricted. It is not a claim
that $K$ is a generic coarsest bisimulation or a full-abstraction theorem.
It states that the declared finite process factors through $K$, and that the
additional detectability clause is sufficient to recover closure equality
from the named productive observations.

\subsection{Operational factorization audit}

The assumption can be checked against a proposed application by asking whether,
conditional on belief and active-project state, every operating reward,
generation law, verification outcome, admission probability, duration law,
completion observation, and next-library update depends on the raw library only
through $(F_A,C_A)$. Provenance, cardinality, identity-specific switching cost,
or a requirement for a named parent violates the assumption unless it is added
to the compressed state. A finite module-gated construction satisfies the
assumption when all projects are indexed only by requirements contained in
$C_A$, incumbent rewards use only $F_A$, and all subsequent kernels use only
belief, active-project variables, and those two objects.

The following exact countermodel shows why the checklist is substantive. Let
$s_0$ be inactive and let active strategies $a$ and $b$ have identical zero
profiles and the same single module. Then libraries $\{s_0,a\}$ and
$\{s_0,b\}$ have identical frontier and closure. If a project is available only
when the raw identity $a$ is present, their project menus differ. The model
therefore does not factor through $(F,C)$, and the projection theorem makes no
claim about it. The countermodel changes no theorem: it falsifies an omitted
assumption, not the stated conditional result.
